% ============================================================
%  PRÉAMBULE — TAC-TIC IntelliPM — Rapport PFE 2026
%  Institut Supérieur des Études Technologiques (ISET)
% ============================================================

% --- Encodage et langue ---
\usepackage[T1]{fontenc}
\usepackage[utf8]{inputenc}
\usepackage[french]{babel}
\frenchbsetup{StandardLists=true}

% --- Géométrie de la page ---
\usepackage[top=2.5cm, bottom=2.5cm, left=3cm, right=2.5cm,
            headheight=15pt]{geometry}

% --- Typographie ---
\usepackage{microtype}
\usepackage{lmodern}
\usepackage{setspace}
\onehalfspacing

% --- Mathématiques ---
\usepackage{amsmath, amssymb}

% --- Couleurs ---
\usepackage[table, dvipsnames]{xcolor}
\definecolor{primary}{RGB}{0,84,166}
\definecolor{lightgray}{RGB}{242,242,242}
\definecolor{codebg}{RGB}{248,248,250}
\definecolor{codeframe}{RGB}{210,215,220}
\definecolor{chapcolor}{RGB}{0,84,166}

% --- Tableaux ---
\usepackage{booktabs}
\usepackage{tabularx}
\usepackage{array}
\usepackage{longtable}
\usepackage{multirow}
\usepackage{colortbl}
\usepackage{float}
\usepackage{makecell}
\renewcommand\theadfont{\bfseries\small}
\renewcommand\theadgape{\Gape[4pt]}
\renewcommand{\arraystretch}{1.3}

% Colonnes personnalisées (C = centré, L = gauche avec largeur fixe)
\newcolumntype{C}[1]{>{\centering\arraybackslash}p{#1}}
\newcolumntype{L}[1]{>{\raggedright\arraybackslash}p{#1}}

% --- Figures ---
\usepackage{graphicx}
\usepackage[font=small, labelfont=bf, labelsep=period]{caption}
\usepackage{subcaption}
\graphicspath{{figures/}}

% --- Boîtes et encadrés ---
\usepackage[framemethod=TikZ]{mdframed}

% --- Listes ---
\usepackage{enumitem}
\usepackage{pifont}
\newcommand{\cmark}{\textcolor{green!60!black}{\ding{51}}}
\newcommand{\xmark}{\textcolor{red!80!black}{\ding{55}}}

% --- Code source ---
\usepackage{listings}

\lstset{
  basicstyle    = \ttfamily\small,
  backgroundcolor = \color{codebg},
  frame         = single,
  rulecolor     = \color{codeframe},
  numbers       = left,
  numberstyle   = \tiny\color{gray!70},
  numbersep     = 8pt,
  breaklines    = true,
  breakatwhitespace = false,
  captionpos    = b,
  showstringspaces = false,
  tabsize       = 2,
  xleftmargin   = 14pt,
  xrightmargin  = 4pt,
  aboveskip     = 6pt,
  belowskip     = 6pt,
}

\lstdefinestyle{sqlstyle}{
  language      = SQL,
  keywordstyle  = \color{blue!80!black}\bfseries,
  commentstyle  = \color{gray}\itshape,
  stringstyle   = \color{orange!70!black},
  morecomment   = [l]{--},
}

\lstdefinestyle{jsstyle}{
  language      = JavaScript,
  keywordstyle  = \color{blue!80!black}\bfseries,
  commentstyle  = \color{gray}\itshape,
  stringstyle   = \color{orange!70!black},
  morecomment   = [l]{//},
}

\lstdefinestyle{phpstyle}{
  language      = PHP,
  keywordstyle  = \color{blue!80!black}\bfseries,
  commentstyle  = \color{gray}\itshape,
  stringstyle   = \color{orange!70!black},
  morecomment   = [l]{//},
  morestring    = [b]{"},
  morestring    = [b]{'},
}

% --- Bibliographie ---
\usepackage[backend=biber, style=ieee, sorting=none, maxbibnames=5]{biblatex}
\addbibresource{references.bib}

% --- En-têtes et pieds de page ---
\usepackage{fancyhdr}
\pagestyle{fancy}
\fancyhf{}
\fancyhead[L]{\small\nouppercase{\leftmark}}
\fancyhead[R]{\small\textit{TAC-TIC IntelliPM — PFE 2026}}
\fancyfoot[C]{\thepage}
\renewcommand{\headrulewidth}{0.4pt}
\renewcommand{\footrulewidth}{0pt}

\fancypagestyle{plain}{%
  \fancyhf{}%
  \fancyfoot[C]{\thepage}%
  \renewcommand{\headrulewidth}{0pt}%
}

% --- Liens hypertexte ---
\usepackage[hidelinks,
            pdftitle={TAC-TIC IntelliPM — Rapport de Projet de Fin d'Études},
            pdfauthor={Fatma Rejeb},
            pdfsubject={Plateforme de gestion de projets intelligente}]{hyperref}

% --- Numérotation ---
\setcounter{tocdepth}{3}
\setcounter{secnumdepth}{3}

% --- Style des titres de chapitre ---
\usepackage{titlesec}

\titleformat{\chapter}[display]
  {\normalfont\huge\bfseries\color{chapcolor}}
  {\filleft\large\MakeUppercase{\chaptertitlename}\ \thechapter}
  {8pt}
  {\filright\Huge}
  [\vspace{4pt}\color{chapcolor}\rule{\linewidth}{1.2pt}\vspace{-2pt}]

\titlespacing*{\chapter}{0pt}{-20pt}{28pt}

\titleformat{\section}
  {\normalfont\large\bfseries\color{chapcolor!85!black}}
  {\thesection}{1em}{}

\titleformat{\subsection}
  {\normalfont\normalsize\bfseries}
  {\thesubsection}{1em}{}

\titleformat{\subsubsection}
  {\normalfont\normalsize\itshape\bfseries}
  {\thesubsubsection}{1em}{}

% --- Environnements personnalisés ---

% Encadré d'introduction de chapitre (fond bleu clair)
\newenvironment{chapintro}{%
  \vspace{4pt}%
  \begin{mdframed}[
    linecolor   = primary!40,
    linewidth   = 0.8pt,
    backgroundcolor = primary!5,
    leftmargin  = 0pt,
    rightmargin = 0pt,
    innerleftmargin  = 14pt,
    innerrightmargin = 14pt,
    innertopmargin   = 10pt,
    innerbottommargin= 10pt,
    skipabove   = 6pt,
    skipbelow   = 8pt,
  ]%
  \itshape
}{%
  \end{mdframed}%
  \vspace{6pt}%
}

% Encadré de conclusion de chapitre (fond bleu, titre "Conclusion du chapitre")
\newenvironment{chapconclusion}{%
  \vspace{10pt}%
  \begin{mdframed}[
    linecolor   = primary,
    linewidth   = 1.5pt,
    backgroundcolor = primary!7,
    leftmargin  = 0pt,
    rightmargin = 0pt,
    innerleftmargin  = 14pt,
    innerrightmargin = 14pt,
    innertopmargin   = 12pt,
    innerbottommargin= 12pt,
    skipabove   = 8pt,
    skipbelow   = 8pt,
    frametitle  = {\textbf{\textcolor{primary}{Conclusion du chapitre}}},
    frametitleaboveskip = 8pt,
    frametitlebelowskip = 6pt,
  ]%
}{%
  \end{mdframed}%
}

% Encadré de problématique (ligne bleue épaisse)
\newmdenv[
  linecolor   = primary,
  linewidth   = 1.5pt,
  frametitle  = {},
  skipabove   = 8pt,
  skipbelow   = 8pt,
]{problembox}

% --- Macro utilitaire : placeholder image manquante ---
% Utilisé dans les chapitres comme \imgplaceholder{largeur}{fichier}{légende}
\newcommand{\imgplaceholder}[3]{%
  \fbox{\parbox{#1}{\centering\vspace{1.2cm}
    \textbf{[IMAGE À INSÉRER]}\\[0.4em]
    \texttt{#2}\\[0.3em]
    \textit{#3}
  \vspace{1.2cm}}}%
}
